
\frontslide{Rappels}



\begin{frame}{Vue d'ensemble d'un système relationnel}


\figSlideWithSize{physique-logique-sql.pdf}{12cm}

\vfill


\begin{blocgauche}
\alert{Retenir}\,: indépendance entre le niveau logique et le niveau physique. Condition
nécessaire de l'administration des SGBD.
\end{blocgauche}

\end{frame}

\begin{frame}{Fonctionnalités d'un SGBD}

\alert{Stockage}\,: les données sont stockées de manière ordonnée sur des supports persistants

\vfill

\alert{Indexation}\,: des structures de données permettent des accès très rapides à la base

\vfill

\alert{Interrogation et mise à jour}\,: les requêtes de haut niveau d'abstraction (SQL) sont transcrites en algorithmes 
efficaces

\vfill

\alert{Partage / Concurrence d'accès}\,: de très nombreuses applications peuvent accéder simultanément à une même base


\vfill
\begin{blocgauche}
\alert{Sécurité}\,: contrôle des accès et gestion des pannes logicielles ou matérielles
\end{blocgauche}

\end{frame}

\begin{frame}[fragile]{Rappels sur le modèle relationnel (par l'exemple)}

Les données sont structurées en \alert{tables} (relations) et \alert{lignes} (nuplets). 
La ligne est l'unité d'information.

\vfill

\begin{small}
\begin{tabular}{l|l|l|c}
\textbf{idArtiste} &  nom & prénom &  annéeNaiss \\
\hline
100 & Eastwood & Clint  & 1930 \\
101 & Hackman & Gene  & 1930 \\
102 & Scott & Tony  & 1930 \\
103 & Smith  & Will  & 1968 \\
\hline
\end{tabular}
\end{small}

\vfill
Les données sont structurées et contraintes par un \alert{schéma}.

\vfill

\begin{minted}[baselinestretch=.9,
               fontsize=\small,
               framesep=2mm]{sql}
create table Artiste  (idArtiste INTEGER NOT NULL,
                   nom VARCHAR (30) NOT NULL,
                   prénom VARCHAR (30) NOT NULL,
                   annéeNaissance INTEGER,
                   PRIMARY KEY (idArtiste));
\end{minted}

\end{frame}


\begin{frame}[fragile]{Clé primaire, clé étrangère}

Les tables sont liées par des valeurs partagées entre clé primaire
et clé étrangère.

\vfill

\begin{small}
\begin{tabular}{l|l|c|l|c|c}
\textbf{idFilm} & titre & année  & genre & \textit{idMES} & \textit{codePays}\\
\hline
20 & Impitoyable & 1992 & Western & 100 & USA \\
21 & Ennemi d'état & 1998 & Action & 102 & USA\\
\hline
\end{tabular}
\end{small}
\vfill


\vfill

\begin{minted}[baselinestretch=.9,
               fontsize=\small,
               framesep=2mm]{sql}
create table Film  (idFilm integer NOT NULL,
                    titre    varchar (80) NOT NULL,
                    année    integer NOT NULL,
                    genre varchar (20) NOT NULL,
                    idRéalisateur    integer,
                    codePays    varchar (4),
                    primary key (idFilm),
                    foreign key (idRéalisateur) 
                      references Artiste(idArtiste),
                    foreign key (codePays) 
                      references Pays(code));

\end{minted}

\end{frame}

\begin{frame}{Les langages}

Le modèle relationnel s'appuie sur deux langages très différents mais \alert{équivalents} 
en pouvoir expressif.

\vfill

\begin{redbullet}

   \item SQL a pour fondement la logique des prédicats. 
   
       \medskip
       
       C'est un langage \alert{déclaratif} qui indique ce que l'on veut obtenir
          sans spécifier \alert{comment} on calcule le résultat.
          
          
        \item L'\alert{algèbre relationnel} est un langage fonctionnel composé
            d'opérateurs. 
            
            
            L'algèbre permet de construire des \alert{plans d'exécutions} indiquant
            comment on évalue une requête.
\end{redbullet}


\vfill
\begin{blocgauche}

\begin{defblock}{Evaluation et optimisation}
Le SGBD transforme une requête au niveau logique (SQL) en un programme d'évaluation
algébrique au niveau physique.
\end{defblock}


\end{blocgauche}

\end{frame}


\begin{frame}[fragile]{Le langage SQL: supposé connu\,!}

Une sélection
\vfill

\begin{minted}[baselinestretch=.9,
               fontsize=\small,
               framesep=2mm]{sql}
select titre
from Film
where année = 2016
\end{minted}

\vfill

Une jointure

\vfill

\begin{minted}[baselinestretch=.9,
               fontsize=\small,
               framesep=2mm]{sql}
select titre, prénom, nom
from Film as f, Artiste as a
where f.idRéalisateur = a.idArtiste
and  année = 2016
\end{minted}

\vfill

La méthode de calcul n'est jamais spécifiée.

\end{frame}


\frame{\frametitle{Algèbre relationnelle\,: : supposé connue également}

Opérateurs qui prennent une ou deux tables en entrée et produisent une table en sortie.

\vfill

\begin{redbullet}
	\item Sélection($\sigma$) : Restriction sur valeur des tuples
	\item Projection($\pi$) : Suppression de colonnes (allège le résultat)
	\item Jointure($\Join$) : Fusionne deux ensembles sur valeur
	\item Différence($-$) : Sous ensemble non présent dans le second
	\item Union($\cup$) : Union des deux ensembles
\end{redbullet}

\vfill

Les opérateurs sont \alert{composables} pour construire des \alert{expressions}
\begin{blocgauche}
$$\pi_{titre} (\sigma_{ann\acute{e}e = 2016} (Film))$$
$$\pi_{titre, pr\acute{e}nom, nom} (\sigma_{ann\acute{e}e = 2016} (Film) \Join_{idR\acute{e}alisateur=idArtiste} Artiste)$$
\end{blocgauche}
}



\begin{frame}{À retenir}

Suivre et comprendre un cours sur les aspects systèmes des SGBDs relationnels suppose
des connaissances préalables


\vfill

\begin{redbullet}
  \item Sur l'architecture des systèmes relationnels, et notamment \alert{l'indépendance entre niveau physique et niveau logique}
  
   \item Sur les principes de conception des schémas relationnels et leurs contraintes
   
   \item Le langage SQL
   \item Et l'algèbre relationnelle
\end{redbullet}

\vfill

À réviser d'urgence si vous avez des doutes ou des lacunes.
\end{frame}
